\documentclass[11pt]{article}

\usepackage[margin=1in]{geometry}
\usepackage{amsmath,amssymb}
\usepackage{booktabs}
\usepackage[T1]{fontenc}
\usepackage[utf8]{inputenc}
\usepackage{lmodern}
\usepackage{microtype}
\usepackage{hyperref}
\usepackage{enumitem}
\usepackage{tabularx}

\hypersetup{
  colorlinks=true,
  linkcolor=blue,
  urlcolor=blue
}

\title{Current Results and Execution Status of
  \texttt{scripts-paper/run\_pipeline.R}}
\author{}
\date{18 July 2026}

\begin{document}

\maketitle

\begin{abstract}
The current paper pipeline does not complete, so it has no current
positive-slack identified-set result, log-variance result, dynamics result, or
validated final artifact route. This is not a matter of missing files in the
report. The current numerical policy deliberately marks every
positive-slack mean-equation profile as \texttt{unreliable}, because the
available local solver cannot certify a global optimum for the indefinite
quadratic problem. Its tau-star search now reports the typed status
\texttt{unresolved}, with certified lower floor zero and no certified upper
transition. The runner then stops while building the structural-table caption:
that caption requests a removed scalar tau-star field, and the typed formatter
receives \texttt{NULL}.

An isolated run of the current source did reach the upstream mean calculation.
With quarterly log consumption growth, the mean sample has $N=256$ quarters,
1962~Q1--2025~Q4. The closed-form $\tau=0$ news-loading point is approximately
$(0.0156,0.0073,-0.4314)$. All three news-loading rows at every displayed
positive slack are \texttt{unreliable}; no positive-slack endpoint
calibration is available. All 200 bootstrap tau-star statuses were
\texttt{unresolved}, with lower-floor band $[0,0]$. Structural tables,
variance shares, log-OLS, PPML, and every later stage are not reached, so no
later empirical headline is a current pipeline result.
\end{abstract}

\tableofcontents

\section{Provenance and the meaning of ``current''}

This note reports an isolated execution of the shared working tree inspected
on 18 July 2026. The tree contains changes beyond \texttt{HEAD} commit
\texttt{cb7c8c6cad920db3fc6d718f9439fbb09f7f3465}. Its current topology is 297
R files: 194 reachable production modules, 100 test modules, and three
inactive or test-support modules; topology checks pass against the 63-record
artifact manifest. In particular, consumption growth now uses
\[
  100\{\log(P_t)-\log(P_{t-1})\},
\]
not $100(P_t/P_{t-1}-1)$, and positive-slack profile candidates now require a
global certificate that the local backend cannot provide.

The verification command was:
\begin{verbatim}
HETID_BOOT_REPS=200 HETID_BOOT_CORES=2 HETID_ASSERT_EQUIV=1 \
  Rscript scripts-paper/run_pipeline.R
\end{verbatim}

The run was performed outside the main working tree so it could not modify
the repository's tracked or ignored paper outputs. It reproduced the current
hard stop described in Section~\ref{sec:failure}.

\section{Execution summary}

\begin{table}[ht]
\centering
\small
\begin{tabularx}{\textwidth}{@{}p{0.29\textwidth}p{0.18\textwidth}X@{}}
\toprule
Stage & Current status & What is reportable \\
\midrule
Data preparation & Completed before halt &
  Quarterly ACM reuse, fitted SDF proxies, consumption log growth, yield
  volatility, and the three PC families were constructed. \\
Mean OLS and $\tau=0$ system & Completed before halt &
  OLS diagnostics, the closed-form $\tau=0$ point, and its point-bootstrap
  intervals are available as pre-halt diagnostics. \\
Positive-slack mean profiles & Unavailable &
  Every news-coefficient row is \texttt{unreliable}; local candidate values
  are deliberately withheld from tables. \\
Positive-slack mean inference & Unavailable &
  No finite certified news-endpoint draws exist, so no endpoint calibration
  is reported. \\
Tau-star status & Completed diagnostic &
  The full sample and all 200 draws are \texttt{unresolved}; their certified
  lower floor is zero, with no certified upper transition. \\
Structural-table caption & Hard failure &
  It passes the removed scalar tau-star field to the typed formatter, which
  receives \texttt{NULL}. \\
Variance shares, mean regions, log-OLS, PPML, and Harvey &
  Not reached &
  No current-run findings from these stages are available. \\
LAD, joint diagnostics, dynamics, EGARCH, log-variance bootstrap &
  Not reached &
  No current-run findings or route claims are available. \\
Final route validation & Not reached &
  No current \texttt{conditional\_route\_status.rds} was written. \\
\bottomrule
\end{tabularx}
\caption{What the current runner did and did not produce.}
\end{table}

\section{Current scientific contract}

The mean equation is
\[
  \Delta c_{t+1}
  =b_0+PC_{E,t}'b_E+PC_{N,t+1}'b_N+\varepsilon_{t+1}.
\]
$PC_E$ comes from a lagged, in-sample fitted expected-SDF proxy, and $PC_N$
from a centered second-order SDF-news proxy. They are generated regressors,
not exact observed conditional expectations or revisions. The specification
maintains $\beta_2^R=0$ for the projection of the news PCs on the specified
common design $(1,PC_E)$. That is an identifying restriction, not a
consequence of unconditional PCA orthogonality.

The single heteroskedasticity instrument is de-meaned five-year realized yield
volatility, \texttt{y60\_vol}. At $\tau=0$, the moment system point-identifies
$b_N$ through a full-rank linear system. At $\tau>0$, the target becomes the
intersection of three indefinite quadratic inequalities. The current local
backend does not certify their coordinate extrema, so every positive-slack
news row fails closed as \texttt{unreliable}.

\section{Results available before the halt}

\subsection{Mean equation}

The mean sample contains 256 quarters, 1962~Q1--2025~Q4. Table
\ref{tab:mean-current} reports only objects that were available before the
failure. The $\tau=0$ intervals are nominal point diagnostics: the closed-form
point plus or minus the standard-normal 95th percentile times the
median-absolute-deviation scale of 200 moving-block point draws. They are not
a joint confidence region and are conditional on the generated SDF panels and
instrument.

\begin{table}[ht]
\centering
\small
\begin{tabular}{lrrrr}
\toprule
Coefficient & OLS & OLS $t$ & $\tau=0$ point & Nominal 90\% point interval \\
\midrule
$b_0$       & 0.787 & 14.41 & 0.787  & $(0.687,\,0.888)$ \\
$b_{1,E}$   & 0.001 & 0.13  & 0.001  & $(-0.007,\,0.009)$ \\
$b_{2,E}$   & -0.461 & -4.52 & -0.475 & $(-0.744,\,-0.206)$ \\
$b_{3,E}$   & 0.903 & 3.07  & 0.810  & $(0.066,\,1.553)$ \\
\midrule
$b_{1,N}$   & 0.0158 & 4.13  & 0.0156 & $(-0.000,\,0.031)$ \\
$b_{2,N}$   & 0.0268 & 0.98  & 0.0073 & $(-0.142,\,0.157)$ \\
$b_{3,N}$   & -0.1820 & -1.70 & -0.4314 & $(-0.756,\,-0.107)$ \\
\bottomrule
\end{tabular}
\caption{Current pre-halt mean-equation diagnostics. OLS has $R^2\approx0.10$
and $N=256$. The $\tau=0$ column is the closed-form identified point. No
positive-slack range is certified.}
\label{tab:mean-current}
\end{table}

The instrument-relevance diagnostics computed before the halt were:
\[
\begin{array}{c|ccc}
 & PC_{N,1} & PC_{N,2} & PC_{N,3}\\
\hline
\operatorname{Cor}(Z,W_{2,i}^2) & 0.567 & 0.484 & 0.275\\
t\text{-statistic in }W_{2,i}^2\text{ on }Z & 10.97 & 8.81 & 4.56\\
\operatorname{Cor}(W_1,W_{2,i}) & 0.168 & 0.035 & -0.062
\end{array}
\]
These describe the aligned sample. They do not turn a local positive-slack
profile into a globally certified bound.

\subsection{What is not a mean-equation result}

The local solver internally found finite baseline candidates, but current code
correctly labels them \texttt{unreliable}. They are not reported here as
identified-set ranges. In particular:

\begin{itemize}[leftmargin=1.5em]
  \item there is no certified statement that any positive-slack news-loading
    range excludes zero;
  \item there is no positive-slack endpoint calibration;
  \item there is no certified news-block or news-component variance-share
    range;
  \item there is no certified pairwise or three-dimensional region; and
  \item there is no estimated boundedness transition $\widehat\tau^\ast$.
\end{itemize}

The tau-star routine now returns
\[
  \texttt{status=unresolved},\qquad
  \texttt{lower=0},\qquad
  \texttt{upper=NA}.
\]
All 200 moving-block draws had the same status and lower floor, so the reported
lower-floor band is $[0,0]$ and zero percent reach the $0.05$ baseline. This
certifies only the $\tau=0$ point. It supplies neither a transition estimate
nor evidence about positive-slack boundedness.

\subsection{No current log-variance diagnostic}

The final live runner stops before loading log-OLS. It therefore produces no
current log-variance sample, point column, residual-crossing census, or mapped
positive-slack range. Log-OLS numbers generated by an immediately preceding
snapshot are historical diagnostics, not results of the source state
documented here.

\section{Current hard failure}
\label{sec:failure}

The failure occurs at runner line~38, immediately after the mean bootstrap,
while the structural-table caption is being built:
\begin{quote}
\path{run_pipeline.R:38}\\
$\longrightarrow$
\path{mean_equation/tables/render_structural_equation_table.R}\\
$\longrightarrow$
\path{mean_equation/tables/structural_equation_caption.R}\\
$\longrightarrow$ \texttt{tau\_transition\_range\_label()}.
\end{quote}

The estimated result stores \texttt{tau\_star\_status},
\texttt{tau\_star\_lower}, and \texttt{tau\_star\_upper}. The caption instead
passes the removed \texttt{set\_id\_mean\_eq\$tau\_star} field to the typed
formatter. The absent field evaluates to \texttt{NULL}, and the formatter's
\texttt{switch()} stops with:

\begin{quote}
\texttt{EXPR must be a length 1 vector}
\end{quote}

The correct tau-star statement is ``unresolved, certified only through
$\tau=0$.'' The current control flow cannot render that statement because the
producer and caption use different object shapes.

A second, later incompatibility remains in PPML. An immediately preceding
snapshot reached the fail-closed log-variance engine, which returned completed
NA/\texttt{unreliable} records without selector diagnostics. PPML coverage
then stopped with
\texttt{coverage selector provenance is absent from a completed engine run}.
That error remains relevant source evidence, but the final live runner does
not reach it.

\section{Downstream findings are unavailable}

Because the runner stops in structural-table rendering, the following old
numerical claims are not current results:

\begin{itemize}[leftmargin=1.5em]
  \item completed structural-table output or variance-share ranges;
  \item log-OLS coefficients or residual-zero crossing counts;
  \item PPML or Harvey attained ranges, standard errors, or bootstrap
    envelopes;
  \item LAD ranges or a claim that the LAD branch ran;
  \item joint-null distances or joint-GMM replication residuals;
  \item Ljung--Box dynamics $p$-values or an EGARCH routing verdict;
  \item log-variance block-length sensitivities;
  \item fitted-volatility envelopes; and
  \item a final conditional route state.
\end{itemize}

The live joint-null code would write an
\texttt{unavailable/not\_evaluated\_uncertified\_boxes} record if reached.
The live region and fitted-volatility code would render ``Unavailable''
placeholders. Those are intended future control-flow outcomes, not artifacts
from the current failed run.

\section{Artifact status}

The manifest still declares 63 paths: 53 required, six conditional on LAD,
and four conditional on EGARCH. The current execution does not validate a
53-file or 59-file materialization. Startup removes all ten conditional paths,
and the hard stop occurs before
\path{output/state/conditional_route_status.rds} is written.

Tracked TeX and SVG outputs in \path{scripts-paper/output/} were produced by
an older workflow that used arithmetic consumption growth and promoted local
positive-slack candidates. Ignored CSV, RDS, and PDF files may also predate
the current source. They must not be used as evidence for current empirical
claims merely because they exist on disk.

\begin{table}[ht]
\centering
\small
\begin{tabularx}{\textwidth}{@{}p{0.34\textwidth}X@{}}
\toprule
Existing artifact family & Current interpretation \\
\midrule
\texttt{structural\_eq*.tex} &
  Historical if it contains positive-slack numbers or old arithmetic-growth
  values. The live renderer is intended to show \texttt{unreliable}, but its
  caption currently fails before the table is completed. \\
\texttt{var\_share*.tex} &
  Historical if it reports positive-slack share ranges; current inputs are
  uncertified. \\
\texttt{log\_var\_eq*}, Harvey, LAD, and inference panels &
  Historical and not regenerated by a completed current run. \\
Descriptive reports &
  Consumption-dependent rows are stale until regenerated from log growth by a
  complete or separately validated current run. \\
Variance-bound summary &
  Its formula is independent of consumption growth, but it was not reached in
  the current run; moreover, it is a leading plug-in approximation, not a
  literal finite-sample upper bound on an omitted remainder. \\
\bottomrule
\end{tabularx}
\caption{How to interpret files left by earlier executions.}
\end{table}

\section{Publication conclusion}

There is no defensible current positive-slack or log-variance headline. The
only structural result presently certified by the code is the closed-form
$\tau=0$ mean point, and even its nominal point intervals remain conditional
diagnostics rather than a joint confidence statement. The current OLS and
$\tau=0$ numbers in Table~\ref{tab:mean-current} may be used to diagnose the
new log-growth specification, but they should not be presented as the output
of a successful end-to-end paper pipeline.

Before publication, three distinct issues must be resolved:

\begin{enumerate}[leftmargin=1.5em]
  \item make the structural caption consume the typed tau-star status and
    endpoints, rather than the removed scalar field;
  \item supply a rigorous global certificate for any positive-slack bound that
    is to be called an identified-set endpoint, or continue reporting those
    objects as unavailable; and
  \item make PPML coverage accept the selector-free fail-closed engine state,
    so the runner can reach truthful downstream unavailable handlers and final
    artifact validation.
\end{enumerate}

After that change, the full current-source workflow must be rerun and all
required and conditional artifacts regenerated before PPML, Harvey, LAD,
joint, dynamics, routing, sensitivity, or region statements can re-enter a
results document.

\section*{References}

\begin{itemize}[leftmargin=1.5em]
  \item Imbens, G. W., and Manski, C. F. (2004). Confidence intervals for
    partially identified parameters. \emph{Econometrica}, 72(6), 1845--1857.
  \item Lewbel, A. (2012). Using heteroscedasticity to identify and estimate
    mismeasured and endogenous regressor models.
    \emph{Journal of Business \& Economic Statistics}, 30(1), 67--80.
\end{itemize}

\end{document}
